% Part: incompleteness
% Chapter: incompleteness-provability
% Section: fixed-point-lemma

\documentclass[../../../include/open-logic-section]{subfiles}

\begin{document}

\olfileid[hi]{inc}{inp}{fix}

\olsection{नियत-बिंदु उपपत्ति}

\begin{explain}
नियत-बिंदु उपपत्ति कहती है कि किसी भी !!{formula}~$!B(x)$ के लिए कोई
!!a{sentence}~$!A$ ऐसा है कि
$\Th{T} \Proves !A \liff !B(\gn{!A})$, बशर्ते $\Th{T}$,~$\Th{Q}$ का
विस्तार हो। झूठे वाक्य के मामले में हम चाहेंगे कि $!A$,~$\Th{T}$ में
सिद्धतः, ``$\gn{!A}$ असत्य है'' के समतुल्य हो; अर्थात् उस कथन के समतुल्य
कि $\Gn{!A}$ किसी असत्य !!{sentence} की ग्योडेल (G\"odel) संख्या है।
प्रमाण का विचार समझने के लिए इसकी तुलना क्वाइन द्वारा~$!A$ की इस अनौपचारिक
व्याख्या से करना उपयोगी होगा: ```अपना उद्धरण पहले रखने पर असत्य देता है'
अपना उद्धरण पहले रखने पर असत्य देता है।'' किसी व्यंजक को लेकर, उसके पहले
उसी व्यंजक का उद्धरण रखकर वाक्य बनाने की संक्रिया को उस व्यंजक का
\emph{विकर्णीकरण} और परिणाम को उसका विकर्ण रूप कह सकते हैं। अतः `अपना
उद्धरण पहले रखने पर असत्य देता है' का विकर्णीकरण है `` `अपना उद्धरण पहले
रखने पर असत्य देता है' अपना उद्धरण पहले रखने पर असत्य देता है।'' अब ध्यान
दें कि क्वाइन का झूठा वाक्य `असत्य देता है' का नहीं, बल्कि `अपना उद्धरण
पहले रखने पर असत्य देता है' का विकर्णीकरण है। इसलिए जिस गुण का
विकर्णीकरण करके झूठा वाक्य मिलता है, उसमें स्वयं विकर्णीकरण शामिल है!{}

अंकगणित की भाषा में हम एक मुक्त चर वाले !!a{formula} की ग्योडेल (G\"odel)
संख्या संगणित करके और उस ग्योडेल (G\"odel) संख्या का मानक अंक-सूचक मुक्त चर के
स्थान पर रखकर उसका उद्धरण बनाते हैं। $!E(x)$ का विकर्णीकरण
$!E(\num{n})$ है, जहाँ $n = \Gn{!E(x)}$। (अब से
$\num{\Gn{!E(x)}}$ को $\gn{!E(x)}$ लिखें।) अतः यदि $!B(x)$ का अर्थ
``असत्य है'' है, तो ``अपना उद्धरण पहले रखने पर असत्य देता है'' का अर्थ
होगा ``अपने विकर्ण रूप की ग्योडेल (G\"odel) संख्या पर लागू किए जाने पर असत्य देता
है।'' यदि हमारे पास ऐसे फलन $\fn{diag}(n)$ के लिए चिह्न~$\Obj{diag}$ होता
जो ग्योडेल (G\"odel) संख्या~$n$ वाले !!{formula} के विकर्ण रूप की ग्योडेल (G\"odel) संख्या
संगणित करता है, तो हम $!E(x)$ को $!B(\Obj{diag}(x))$ लिख सकते थे। तब झूठे
वाक्य का क्वाइन रूप उसका विकर्णीकरण होता, अर्थात्
$!E(\gn{!E(x)})$ अथवा
$!B(\Obj{diag}(\gn{!B(\Obj{diag}(x))}))$। निस्संदेह अब $!B(x)$ कोई अन्य
गुण भी हो सकता है और यही रचना काम करेगी। अपूर्णता-प्रमेय के लिए हम
$!B(x)$ को ``$x$,~$\Th{T}$ में !!{derivable} नहीं है'' लेंगे। तब $!E(x)$
का अर्थ होगा ``अपने विकर्ण रूप की ग्योडेल (G\"odel) संख्या पर लागू किए जाने पर ऐसा
!!a{sentence} देता है जो~$\Th{T}$ में !!{derivable} नहीं है।''

इसे~$\Th{T}$ में औपचारिक बनाने के लिए हमें $\fn{diag}$ को औपचारिक बनाने
का ढंग खोजना होगा। फलन $\fn{diag}(n)$ संगणनीय, बल्कि आदिम पुनरावर्ती है:
यदि $n$ किसी सूत्र~$!E(x)$ की ग्योडेल (G\"odel) संख्या है, तो $\fn{diag}(n)$,
$!E(\gn{!E(x)})$ की ग्योडेल (G\"odel) संख्या लौटाता है। (याद करें,
$\gn{!E(x)}$,~$!E(x)$ की ग्योडेल (G\"odel) संख्या का मानक अंक-सूचक, अर्थात्
$\num{\Gn{!E(x)}}$, है।) यदि $\Obj{diag}$,~$\Th{T}$ में फलन
$\fn{diag}$ का निरूपण करने वाला फलन-चिह्न होता, तो हम $!A$ को सूत्र
$!B(\Obj{diag}(\gn{!B(\Obj{diag}(x))}))$ ले सकते थे। ध्यान दें कि
\begin{align*}
\fn{diag}(\Gn{!B(\Obj{diag}(x))}) & =
\Gn{!B(\Obj{diag}(\gn{!B(\Obj{diag}(x))}))} \\
& = \Gn{!A}.
\end{align*}
यह मानते हुए कि $\Th{T}$
\[
\Obj{diag}(\gn{!B(\Obj{diag}(x))}) = \gn{!A},
\]
को !!{derive} कर सकता है, वह
$!B(\Obj{diag}(\gn{!B(\Obj{diag}(x))}))
\liff !B(\gn{!A})$ को !!{derive} कर सकता है। किंतु बायाँ पक्ष परिभाषा से
$!A$ है।

निस्संदेह $\Obj{diag}$ सामान्यतः $\Th{T}$ का फलन-चिह्न नहीं होगा, और
निश्चिततः वह~$\Th{Q}$ का फलन-चिह्न नहीं है। किंतु $\fn{diag}$ संगणनीय है,
इसलिए उसका~$\Th{Q}$ में किसी सूत्र $!D_{\fn{diag}}(x,y)$ द्वारा
\emph{निरूपण} होता है। अतः $!B(\Obj{diag}(x))$ लिखने के स्थान पर हम
$\lexists[y][(!D_{\fn{diag}}(x,y) \land !B(y))]$ लिख सकते हैं। ऊपर
रेखांकित शेष प्रमाण फिर चल जाता है, और वास्तव में वह~$\Th{Q}$ में ही चल
जाता है।
\end{explain}

\begin{lem}
\ollabel{lem:fixed-point} $!B(x)$ एक मुक्त चर~$x$ वाला कोई सूत्र हो। तब
कोई वाक्य~$!A$ ऐसा है कि $\Th{Q} \Proves !A \liff !B(\gn{!A})$।
\end{lem}

\begin{proof}
$!B(x)$ दिए जाने पर $!E(x)$ को सूत्र
$\lexists[y][(!D_{\fn{diag}}(x,y) \land !B(y))]$ मानें और $!A$ को उसका
विकर्णीकरण, अर्थात् सूत्र $!E(\gn{!E(x)})$, मानें।

चूँकि $!D_{\fn{diag}}$,~$\fn{diag}$ का निरूपण करता है और
$\fn{diag}(\Gn{!E(x)}) = \Gn{!A}$, इसलिए $\Th{Q}$ निम्न को !!{derive}
कर सकता है:
\begin{align}
  & {}!D_{\fn{diag}}(\gn{!E(x)}, \gn{!A}) \ollabel{repdiag1} \\
  & \lforall[y][(!D_{\fn{diag}}(\gn{!E(x)},y) \lif
  \eq[y][\gn{!A}])]. \ollabel{repdiag2}
\end{align}
अब हम दिखाते हैं कि $\Th{Q} \Proves !A \liff !B(\gn{!A})$। हम केवल तर्क
और~$\Th{Q}$ में !!{derivable} तथ्यों का उपयोग करके अनौपचारिक तर्क देते हैं।

पहले $!A$, अर्थात् $!E(\gn{!E(x)})$, मान लें। $!E(x)$ की परिभाषा पर
वापस जाने से देखते हैं कि $!E(\gn{!E(x)})$ केवल
\[
\lexists[y][(!D_{\fn{diag}}(\gn{!E(x)},y) \land !B(y))].
\]
है। ऐसा~$y$ लें। चूँकि $!D_{\fn{diag}}(\gn{!E(x)},y)$, इसलिए
\olref{repdiag2} से $y = \gn{!A}$। अतः $!B(y)$ से हमें~$!B(\gn{!A})$
मिलता है।

अब $!B(\gn{!A})$ मान लें। \olref{repdiag1} से हमारे पास
\begin{align*}
& {}!D_{\fn{diag}}(\gn{!E(x)}, \gn{!A}) \land !B(\gn{!A}).
\intertext{इससे निकलता है कि}
& \lexists[y][(!D_{\fn{diag}}(\gn{!E(x)},y) \land !B(y))].
\end{align*}
किंतु यह केवल $!E(\gn{!E(x)})$, अर्थात्~$!A$, है।
\end{proof}

\begin{digress}
इसकी तुलना संगणनीयता-सिद्धांत में नियत-बिंदु उपपत्ति के प्रमाण से करें।
अंतर यह है कि यहाँ हम किसी \emph{कथन} को स्वयं उसके पदों में परिभाषित करना
चाहते हैं, जबकि वहाँ हम किसी \emph{फलन} को स्वयं उसके पदों में परिभाषित
करना चाहते थे; इस अंतर के अतिरिक्त विचार वास्तव में एक ही है।
\end{digress}

\begin{prob}
  कोई !!^a{formula}~$!A(x)$ \emph{सत्यता-परिभाषा} है यदि सभी
  !!{sentence}s~$!B$ के लिए
  $\Th{Q} \Proves !B \liff !A(\gn{!B})$। नियत-बिंदु उपपत्ति का उपयोग करके
  दिखाएँ कि कोई !!{formula} सत्यता-परिभाषा नहीं है।
\end{prob}

\end{document}
